\documentclass[11pt]{article}

\usepackage[margin=1in]{geometry}
\usepackage[T1]{fontenc}
\usepackage[utf8]{inputenc}
\usepackage{lmodern}
\usepackage{microtype}
\usepackage{amsmath, amssymb, amsthm}
\usepackage{mathtools}
\usepackage{xcolor}
\usepackage{hyperref}
\usepackage{parskip}

\definecolor{linknavy}{HTML}{1A4FA0}
\hypersetup{
  colorlinks=true,
  urlcolor=linknavy,
  pdftitle={Seminar Report -- Polynomials and High-Dimensional Spheres},
  pdfauthor={Vandan Pokal}
}

\title{Seminar Report \\[2pt]
       \large Polynomials and High-Dimensional Spheres \\
       (talk by Prof.\ Ambar Sengupta)}
\author{Vandan Pokal}
\date{December 7, 2024}

\begin{document}

\maketitle

\section*{Introduction}

This report is on the talk \emph{Polynomials and High-Dimensional Spheres} given
by Prof.\ Ambar Sengupta of the University of Connecticut on November 12, 2024.

The talk opened with a brief historical primer and was organised into four
parts: Gaussian integration and spherical integration; harmonic polynomials
and operators on polynomial spaces; Hermite limits over large spheres; and the
spherical Laplacian in higher dimensions.

Three key results were highlighted. First, the spherical Laplacian converges to
the Hermite differential operator:
\[
  \lim_{N\to\infty} \Delta_{S^{N-1}(\sqrt{N})}
  \;=\; \sum_{i=1}^{\infty}\bigl(\partial_i^{\,2} - X_i\,\partial_i\bigr).
\]
Second, Gegenbauer polynomials converge to Hermite polynomials:
\[
  \lim_{N\to\infty} q_m(\sqrt{N};X)
  \;=\; \lim_{N\to\infty} C_m^{(N-2)/2}\!\bigl(X/\sqrt{N}\bigr)
  \;=\; H_m(X),
\]
where $q_m(\sqrt{N}; X)$ coincides with a zonal harmonic polynomial on the
sphere $S^{N-1}(\sqrt{N})$. Third, Gaussian integration arises as the limit of
spherical integration:
\[
  \langle p, q\rangle_{L^2(\mathbb{R}^\infty,\,\mu)}
  \;=\; \lim_{N\to\infty}
  \langle p, q\rangle_{L^2(S^{N-1}(\sqrt{N}),\,\bar{\sigma})},
\]
where $p, q \in \mathcal{P}_k$ are viewed in $\mathbb{R}^N$.

This third result is remarkable: it simplifies the infinite-dimensional
spherical Laplacian --- which is already cumbersome in finite dimensions --- to
a tractable and concise Hermite differential operator. The use of rotation
generators, the quadratic Casimir, and the restriction--projection map onto the
sphere (which produces a unique harmonic polynomial) was particularly
intuitive. One motivation mentioned in the talk lies in statistical mechanics,
where the total kinetic energy of a system can be written as a function of
velocities indexed by particles. As the number of particles grows, the total
kinetic energy becomes increasingly difficult to handle directly, and the
results of this talk offer a way to simplify these kinds of problems.

The report below summarises each part of the talk. It draws on a careful
reading of the accompanying paper by Amy Peterson and Ambar Sengupta of the
University of Connecticut.

\section*{Gaussian Integration as the Limit of Spherical Integration}

The plan of attack in this part was first to establish a relation between the
spherical integral of a homogeneous (Borel-measurable) polynomial function and
the corresponding Gaussian integral. The relation is intuitive: the Gaussian
integral over $\mathbb{R}^N$ can be written as a coefficient times the integral
of the same function over the unit sphere. Explicitly, the coefficient is
\[
  A_{d,N} \;=\; 2^d \,\frac{\Gamma\!\left(\frac{d+N}{2}\right)}{\Gamma\!\left(\frac{N}{2}\right)},
\]
where $d$ is the degree of the function. The $L^2$ pairing on the sphere is
then related to the Gaussian $L^2$ pairing by
\[
  a_{d,N}\,\langle p, q\rangle_{L^2(S^{N-1},\,\bar{\sigma})}
  \;=\; \langle p, q\rangle_{L^2(\mathbb{R}^N,\,\mu)},
\]
where
\[
  a_{d,N} \;=\; \prod_{j=1}^{d}\!\left(1 + \frac{2j-1}{N}\right)
  \qquad\text{and}\qquad d \;=\; \frac{d_p + d_q}{2}.
\]
This can be shown straightforwardly by polar disintegration and simplification
via the Gamma function.

The interesting consequence is that for $p, q \in \mathcal{P}_k$ with $k < N$,
the spherical $L^2$ pairing forms an inner product on $\mathcal{P}_k$. Since
$\mathcal{P}_k = \bigoplus_{d=0}^{\infty} \mathcal{P}^{d}$ and
$a_{d,N} \xrightarrow{N\to\infty} 1$, we obtain
\[
  \lim_{N\to\infty}
  \langle p, q\rangle_{L^2(S^{N-1},\,\bar{\sigma})}
  \;=\; \langle p, q\rangle_{L^2(\mathbb{R}^\infty,\,\mu)}.
\]

\section*{Harmonic Polynomials and Operators on Polynomial Spaces}

This part was dedicated to developing the relevant operators and their actions
on polynomial spaces.

If we endow the polynomial space $\mathcal{P}$ with a Hermite inner product,
then the adjoint of $\partial_j$ is simply multiplication by $X_j$, where $X_j$
is the corresponding variable. It follows easily that the adjoint of $\Delta_N$
is multiplication by $\|X\|^2 = X_1^2 + \cdots + X_N^2$.

Since $\Delta_N : \mathcal{P}_N^{d} \to \mathcal{P}_N^{d-2}$ is a linear
operator, we can split
\[
  \mathcal{P}_N^{d}
  \;=\; \ker \Delta_N \;\oplus\; \|X\|^2 \mathcal{P}_N^{d-2}
  \;=\; \mathcal{H}_N^{d} \;\oplus\; \|X\|^2 \mathcal{P}_N^{d-2}.
\]
This extends naturally to $\mathcal{P}_N^{\le d}$. Thus, for any
$p \in \mathcal{P}_N$, we have $p = p_0 + \|X\|^2 q$ with $p_0 \in \mathcal{H}_N$
and $q \in \mathcal{P}_N$. More generally, any $p \in \mathcal{P}_N^{\le d}$ can
be expanded as
\[
  p \;=\; p_0 \;+\; \|X\|^2 p_1 \;+\; \cdots \;+\; \|X\|^{2s}\, p_s,
\]
where $p_j \in \mathcal{H}_N^{\le d - 2j}$.

The operators $L_a$ (which generates a unique harmonic polynomial), $M_{jk}$
and $\|M\|^2$, together with their respective algebraic properties, were then
developed. These tools were used to derive zonal spherical harmonics. By
definition, a zonal harmonic is a harmonic polynomial that is invariant under
all rotations fixing one vector. To simplify the algebra, we may assume
$t = (1, 0, 0, \ldots, 0)$, since any vector can be obtained from $t$ by scaling
and rotation. If $p \in \mathcal{P}_N$ is invariant under all rotations
preserving a given direction, then $p = q(a; X)$ for all $x \in S^{N-1}(a)$.
On the sphere $L_a p = L_a q(a;X)$, and $L_a q(a;X) \in \mathcal{H}_N$ is
homogeneous. Because this polynomial is invariant under rotations preserving a
direction, harmonic, and homogeneous, it is a zonal harmonic. This is how zonal
harmonics may be generated, with the distinction that when restricted to the
sphere $q(a; X)$ is itself a zonal harmonic. Concretely,
\[
  C_m^{(N-2)/2}\!\bigl(X/a\bigr) \;=\; q_m(a; X)
\]
when restricted via $L_a$. This identity is used later in conjunction with
orthogonality under $\langle\cdot,\cdot\rangle_{S^{N-1}(a)}$ to obtain the
Hermite limit.

Since $\langle\cdot,\cdot\rangle_{a,N}$ is an inner product on $\mathcal{P}_k$
when $k < N$, the adjoints of the operators in use can be found by appealing
to the earlier relation between spherical and Gaussian integration. The
properties of these operators are then used to show that $\mathcal{H}_N^{m}$
is orthogonal to $\mathcal{H}_N^{n}$ when $m \ne n$. From this orthogonality of
the zonal harmonics, one derives $q_m(a; X_1)$, monic with $q_0(X_1) = 1$,
satisfying
\[
  q_m(a; X_1) \;=\; (I - P_{a,N})\, X_1^{m},
\]
where $P_{a,N}$ is the projection from $\mathcal{P}_N^{\le m}$ onto
$\mathcal{P}_N^{\le m - 1}$ with respect to
$\langle\cdot,\cdot\rangle_{a,N}$. Furthermore, applying Gram--Schmidt to
$1, X, X^2, \ldots$ under this inner product produces the zonal harmonics
directly.

\section*{Hermite Limits}

In this part the Gaussian $L^2$ inner product was revisited. The projection
$P_N^{\prime\,\le d} : \mathcal{P}_N^{\le d} \to \mathcal{P}_N^{\le d - 1}$
extends to all of $\mathcal{P}^{\le d}$, since the restriction of the
projection in $M$ variables to $\mathcal{P}_N^{\le d}$ agrees with the
projection in $N$ variables.

Using this, Hermite polynomials are produced from monomials:
\[
  \bigl(1 - P^{\prime\,\le d}\bigr)\bigl(X_1^{m_1}\!\cdots X_N^{m_N}\bigr)
  \;=\; H_{m_1}(X_1) \cdots H_{m_N}(X_N),
\]
with $m_1 + \cdots + m_N = d$.

Using the fact that if a sequence of inner products
$\langle\cdot,\cdot\rangle_n \to \langle\cdot,\cdot\rangle$, then the
orthogonal projections $P_n$ onto a fixed proper subspace converge to the
projection $P$ under the limiting inner product, and combining this with the
limiting relation between the Gaussian and spherical inner products from the
first part, we obtain
\[
  \lim_{N\to\infty}\!\bigl(1 - P_{k,N}^{\le d}\bigr)\bigl(X_1^{m_1} \cdots X_k^{m_k}\bigr)
  \;=\; \lim_{N\to\infty}\!\bigl(1 - P^{\prime\,\le d}\bigr)\bigl(X_1^{m_1} \cdots X_k^{m_k}\bigr)
  \;=\; H_{m_1}(X_1) \cdots H_{m_k}(X_k).
\]
The limiting relation stated in the introduction between zonal harmonics,
Gegenbauer polynomials, and Hermite polynomials then follows.

\section*{Spherical Laplacian}

In this part the properties of the spherical Laplacian were discussed first.
The operator $\Delta_{N-1,a}$ acts on $p \in \mathbb{C}[X_1, \ldots, X_n]$ in
two cases. First, when $n \le N$,
\[
  \Delta_{N-1,a} \;=\; \frac{1}{a^2}\,\|M\|^2,
\]
and second, when $n > N$,
\[
  \Delta_{N-1,a}\, p \;=\; \frac{1}{a^2}\,\|M\|^2 p \;=\; q_{\min},
\]
where $q_{\min} \in \mathbb{C}[X_1, \ldots, X_n]$ is the unique minimal
representative modulo the ideal $\mathcal{Z}_N(a)$ of polynomial multiples of
$\bigl(\|X\|^2 - a^2\bigr)$. This is demonstrated by exploiting the
self-adjointness of the spherical Laplacian with respect to
$\langle\cdot,\cdot\rangle_{a,N}$, together with linearity. Examining the
action on monomials in the limit $N \to \infty$, one finds that
\[
  \Delta_{N-1,\sqrt{N}}
  \;\longrightarrow\;
  \sum_{i=1}^{\infty}\bigl(\partial_i^{\,2} - X_i\,\partial_i\bigr)
\]
for any polynomial in $\mathcal{P}_N^{\le m}$.

\section*{Remarks}

Both the talk and the accompanying paper were straightforward to follow and
genuinely intuitive, thanks to the systematic use of operators on polynomial
spaces.

\section*{References}

Peterson, A., \& Sengupta, A.\,N. (2019).
\emph{Polynomials and High-Dimensional Spheres}.
\href{https://arxiv.org/abs/1903.07697}{arXiv:1903.07697 [math.MP]}.

\end{document}
